% !BIB program = bibtex
%\documentclass[journal,12pn, twocolumn]{IEEEtran}
\documentclass[journal,12pn, onecolumn]{IEEEtran}
\usepackage{array}
\usepackage{multirow}
\usepackage{mathrsfs}
\usepackage{amsmath,amssymb,amsfonts}
\usepackage{caption}
\usepackage{makecell}
\usepackage{float}
\usepackage{textcomp}
\usepackage{cite}
\usepackage{amsthm}
\usepackage{verbatim}
%\usepackage{bbding}
\usepackage{array}
\usepackage{mathrsfs}
%\usepackage{amssymb,amsmath,amsthm}
\usepackage{amsfonts}
\usepackage{booktabs}
\usepackage{graphicx}
\usepackage{epsfig}
\usepackage{diagbox}
\usepackage{amsthm}
\usepackage{multicol}
\usepackage{threeparttable}
\usepackage{algorithm}
\usepackage{algorithmic}
\usepackage{subfigure}
\usepackage{xcolor}
\definecolor{customgreen}{RGB}{0, 128, 0} %
\usepackage{bm}
\usepackage{tabularx}


\usepackage{epstopdf}
\usepackage{url}

\hyphenation{op-tical net-works semi-conduc-tor}
\newtheorem{theorem}{Theorem}
\newtheorem{proposition}{Proposition}
\newtheorem{Definition}{Definition}
\newtheorem{example}{Example}
\renewcommand*{\proofname}{\noindent\bf Proof}
\renewcommand\labelenumi{(\alph{enumi})}
\DeclareMathOperator*{\argmax}{\arg\max}
\DeclareMathOperator*{\argmin}{\arg\min}
\usepackage[colorlinks,linkcolor=orange,anchorcolor=orange,citecolor=orange,]{hyperref}
\hypersetup{hypertex=true,
            colorlinks=true,
            linkcolor=cyan,
            anchorcolor=orange,
            citecolor=green}
\renewcommand{\baselinestretch}{0.96}
\newenvironment{breakablealgorithm}
  {% \begin{breakablealgorithm}
   \begin{center}
     \refstepcounter{algorithm}% New algorithm
     \hrule height.8pt depth0pt \kern2pt% \@fs@pre for \@fs@ruled
     \renewcommand{\caption}[2][\relax]{% Make a new \caption
       {\raggedright\textbf{Solution~\thealgorithm} ##2\par}%
       \ifx\relax##1\relax % #1 is \relax
         \addcontentsline{loa}{algorithm}{\protect\numberline{\thealgorithm}##2}%
       \else % #1 is not \relax
         \addcontentsline{loa}{algorithm}{\protect\numberline{\thealgorithm}##1}%
       \fi
       \kern2pt\hrule\kern2pt
     }
  }{% \end{breakablealgorithm}
     \kern2pt\hrule\relax% \@fs@post for \@fs@ruled
   \end{center}
  }

\newtheorem{remark}{Remark}

\allowdisplaybreaks[4]
\hyphenation{op-tical net-works semi-conduc-tor IEEE-Xplore}

\renewcommand{\algorithmicrequire}{\textbf{Input:}}

\renewcommand{\algorithmicensure}{\textbf{Output:}}
\begin{document}

\newcolumntype{Y}{>{\centering\arraybackslash}X}

% -------- Reply to Editors and Reviewers -------
\title{Responses to the Review Comments of BDR-D-25-00303R3}
% \ \\
% \normalsize Title: \emph{A $K$-means Triangular Synthesis Large Margin Classifier with Unified Pinball Loss for Imbalanced Data}}
% \author[ads]{Danlin Shao}\ead{dlshao1104@163.com}
% \author[ads]{Yixi Dai}\ead{YixiDai2001@163.com}
% \author[ads]{Junjie Li}\ead{JunJieLi45@163.com}
% \author[ads]{Shenglin Li}\ead{lishenglin@swu.edu.cn}
% \author[ads]{Rui Chen}\ead{rchencr@163.com}
% \author{Danlin Shao, Yixi Dai, Junjie Li, Shenglin Li, Rui Chen*}
\maketitle

\textbf{Original Manuscript ID: }BDR-D-25-00303R3 \\


\textbf{Original Article Title: }``Dual Balanced Large Marginal Machine for Imbalanced Data Classification" \\

\textbf{To: }Big Data Research Editor\\

\textbf{Re: }Response to reviewers\\

\section*{}
{
Dear Editor, \\

Thank you for allowing a resubmission of our manuscript, and we particularly appreciate the Editor-in-Chief for granting us the opportunity to conduct a round of revisions to address the reviewers’ comments.\\

We are uploading (a) our point-by-point response to the comments (below) (response to reviewers, under “Author’s Response Files”), (b) an updated manuscript with orange highlighting indicating changes, and (c) an Supplementary file submitted in accordance with the reviewers’ requirements.\\

Best wishes,\\

Pingping Gu, Denghao Dong, Quan Shi, Fudi Wang and Guojun Huang*}\\






\textbf{The response letter is structured as follows:}\\

\begin{itemize}
    \item [$\bullet$]Response for Editor: Page 2.\\
    \item [$\bullet$]Response for Reviewer 7: Page 3.\\
    \item [$\bullet$]Response for Reviewer 8: Page 4-8.\\
\end{itemize}

\newpage

% % --------------- Editor ----------------------

% \section*{\textbf{Section \#1-a: Entire editorial decision letter}}\label{sec:1a}
% Re: \textbf{INS-D-24-1205}

% ~

% Dear Dr. Zhang,

% ~

% We have received the decision on your paper entitled ``A Novel Cost-Sensitive Three-Way Intuitionistic Fuzzy Large Margin Classifier".

% ~

% The Editor's and the reviewers' comments are as follows:


% ~

% The Introduction may be improved to give a reader the context of this study.

% For example,

% (a) in addition to listing a number references, it might be useful to identify a few most relevant studies,

% and (b) the most recent results on the wide sense of three-way decision may be briefly commented.

% ----------

% ~

% Reviewer \#1: This paper introduces three-way decisions into the large margin distribution machine. Generally, the method is interesting, and the results are good. I have a number of comments.

% 1. Three-way decision (3WD) has been found widespread application across diverse fields $\rightarrow$ Three-way decision (3WD) has been widely applied in diverse fields

% 2. costs; iii) the disregard of sample $\rightarrow$ costs; and iii) the disregard of sample

% 3. kaggle $\rightarrow$ Kaggle

% 4. See Section 2.1 $\rightarrow$ see Section 2.1

% 5. The framework of 3WD is shown in Fig. 1. $\rightarrow$ Fig. 1. illustrates the framework of 3WD.

% 6. It is better to present the experimental parameter settings in a table. The reason of respective settings should be explained as well.

% 7. What is the potential real-life applications of the proposed method?

% 8. "transfusion\_new" $\rightarrow$ ``transfusion\_new''. The double quotation marks are incorrect.

% 9. ``breast\_w" is not presented in Table 2.

% 10. The values in tables (e.g., Table 5) should be aligned to the right.

% ~

% Reviewer \#2: I have carefully reviewed the manuscript. The paper proposes an innovative model, the CS3W-IFLMC model, to address uncertainty and noise issues effectively. I think this work is interesting and meaningful, but there are still some parts that require further revision and improvement.

% 1. Please explain how the category centers of the samples are determined in the proposed Intuitionistic fuzzy strategy, which helps readers better understand the construction process of the model.

% 2. Could you please clarify the criteria for using bold font in the table? It should be used to highlight the optimal experimental results. However, I noticed that the number ``89.90" in bold does not represent an optimal result in the experimental findings of the ``kr-vs-kp" dataset in Table 4.

% 3. The description of the experimental results in section 4.2 should be more detailed to highlight the superiority of the proposed IFLMC model.

% 4. Please clarify the meaning of the bold symbol ``X" in Equation 26.

% 5. Some other issues to be corrected.

% 5.1 In the title of Figure 4, ``CMC-0-2" should be corrected to ``cmc-0-2".

% 5.2 A period is missing at the end of the title in Figure 1.

% 5.3 There is a missing space before the fifth sentence of the second paragraph in section 4.3.2.

% 5.4 In section 2.3, ``The 3WD model provide a feasible way..." should be ``The 3WD model provides a feasible way...".

% 5.5 In section 3, ``CS3W-UFLMC" should be ``CS3W-IFLMC".

% ~

% Reviewer \#3: This paper combines the probability output of LDM with the cost-sensitive three-way decision method to achieve delayed decision-making in LDM. Additionally, an innovative dual-center intuitionistic fuzzy method is proposed and introduced, thereby establishing a universal model, CS3W-IFLMC. The paper demonstrates the effectiveness of the methods used both theoretically and experimentally. This study contains solid and original contributions., especially for three-way decision and LDM. I only have a few minor comments as follows.

% 1. To more convincingly demonstrate the superiority of the fuzzy strategy proposed in this paper compared to other strategies, it is necessary to contrast it with existing fuzzy strategies. This comparison will provide a clearer picture of the advantages and disadvantages of each approach, highlighting the unique benefits of the strategy proposed in this paper.

% 2. In the main text, the full name should be used the first time an acronym is introduced. (e.g., introduction, IF set theory)

% 3. A link to the code repository should be included at the end of the article to facilitate readers' access and reproducibility of the model.

% 4. The format of the chart titles should be consistent throughout the entire manuscript. Table 1 follows the convention of title capitalization, while other tables use sentence capitalization.

% 5. Please verify if there are any errors in the formulas. For example, whether ``n" in Formula 5 should be ``m".

% 6. The language should be improved further.


% ~

% \textcolor{red}{Note: Please double check the author names provided in the manuscript source file so that authorship related changes are made in the revision stage. If your manuscript is accepted, any authorship change will involve approval from co-authors and respective editor handling the submission and this may cause a significant delay in publishing your manuscript.}

% ~

% In view of these comments the Editor-in-Chief, Professor Witold Pedrycz, has decided that the paper can be reconsidered for publication after major revisions.

% ~

% When you log-in to your Main Menu, you will see a menu item called 'Submission Needing Revision'. You will find your submission record there. Also, the reviewer(s) may have uploaded a file with detailed comments on your manuscript. Click on "View Reviewer Attachments" to access any detailed comments from the reviewer(s) that may have been included.

% ~

% When submitting your revised manuscript, please ensure that you upload the source files (e.g. Word or. TeX with other files in a.zip file which the system will automatically unzip). Uploading a PDF file at this stage will create delays should your manuscript be finally accepted for publication. If your revised submission does not include the source files, we will contact you to request them.

% ~

% Please note that this journal offers a new, free service called AudioSlides: brief, webcast-style presentations that are shown next to published articles on ScienceDirect (see also \textcolor{orange}{http://www.elsevier.com/audioslides}). If your paper is accepted for publication, you will automatically receive an invitation to create an AudioSlides presentation.

% ~

% The revised version of your submission is due by \textbf{20 Mar 2024}.

% ~

% Information Sciences features the Interactive Plot Viewer, see: \textcolor{orange}{http://www.elsevier.com/interactiveplots}. Interactive Plots provide easy access to the data behind plots. To include one with your article, please prepare a .csv file with your plot data and test it online at \textcolor{orange}{http://authortools.elsevier.com/interactiveplots/verification} before submission as supplementary material.

% ~

% Research Elements (optional)
% This journal encourages you to share research objects - including your raw data, methods, protocols, software, hardware and more – which support your original research article in a Research Elements journal. Research Elements are open access, multidisciplinary, peer-reviewed journals which make the objects associated with your research more discoverable, trustworthy and promote replicability and reproducibility. As open access journals, there may be an Article Publishing Charge if your paper is accepted for publication. Find out more about the Research Elements journals at \textcolor{orange}{https://www.elsevier.com/authors/tools-and-resources/research-elements-journals?dgcid=ec\_em\_research\_elements\_email}.

% ~

% We look forward to receiving the revised version of your paper together with a reply to the reports and a summary of the revisions made.

% ~

% Include interactive data visualizations in your publication and let your readers interact and engage more closely with your research. Follow the instructions here: \textcolor{orange}{https://www.elsevier.com/authors/author-services/data-visualization} to find out about available data visualization options and how to include them with your article.

% ~

% Kind regards,

% ~

% Witold Pedrycz, PhD., DSc

% Editor-in-Chief

% Information Sciences

% ~

% At Elsevier, we want to help all our authors to stay safe when publishing. Please be aware of fraudulent messages requesting money in return for the publication of your paper. If you are publishing open access with Elsevier, bear in mind that we will never request payment before the paper has been accepted. We have prepared some guidelines (\textcolor{orange}{https://www.elsevier.com/connect/authors\\-update/seven-top-tips-on-stopping-apc-scams} ) that you may find helpful, including a short video on Identifying fake acceptance letters (\textcolor{orange}{https://www.youtube.com/watch?v=o5l8thD9XtE} ). Please remember that you can contact Elsevier s Researcher Support team (\textcolor{orange}{https://service.elsevier.com/app/home/supporthub/publishing/}) at any time if you have questions about your manuscript, and you can log into Editorial Manager to check the status of your manuscript (\textcolor{orange}{https://service.elsevier.com/app/answers/detail/a\_id\\/29155/c/10530/supporthub/publishing/kw/status/}).

% \#AU\_INS\#

% ~

% To ensure this email reaches the intended recipient, please do not delete the above code


\newpage


\section*{\textbf{Editor}}\label{edi}

{\color{black}
\subsubsection*{\underline{General Comments}}
\emph{\\
I have completed my evaluation of your manuscript. The reviewers recommend reconsideration of your manuscript following minor revision and modification. I invite you to resubmit your manuscript after addressing the comments below. Please resubmit your revised manuscript by 07/05/2026.
When revising your manuscript, please consider all issues mentioned in the reviewers’ comments carefully: please outline every change made in response to their comments and provide suitable rebuttals for any comments not addressed. Please note that your revised submission may need to be re-reviewed.
}}\\


{\color{black}
\subsubsection*{\textbf{\underline{Response}}}
We sincerely appreciate the Editor-in-Chief and all reviewers for their constructive feedback and for granting us the opportunity to further refine our manuscript. We have carefully addressed all comments raised in this fourth round of review. The detailed point-by-point responses will be added below under Sections ``Reviewer 7'' and ``Reviewer 8''.}\\

{\color{blue}
The point-by-point responses to Reviewer \#7 and Reviewer \#8 will be added in the sections below. Author responses and author actions are left blank for now and will be completed during the revision.
}
{\color{black}
To facilitate reviewers in checking our revisions, we have prepared detailed response documents. The concrete revisions and responses to the reviewers are presented in Sections ``Reviewer 7'' and ``Reviewer 8''. Note that reference numbers are in accordance with the references at the end of \textbf{this response document}. For clarity, we present the section or subsection, page, and line numbers of the revisions based on \textbf{the newly revised version manuscript}, NOT the previous manuscript. Moreover, we mark the revised contents in the manuscript using {\color{magenta} magenta color} to distinguish them from previous revisions. The detailed revisions are presented as follows.
}



\newpage

% ---------------- Reviewer 7 ---------------------

{\color{black}
\section*{\textbf{REVIEWER 7}}\label{rev7}}

{\color{black}
\subsubsection*{\underline{General Comments}}
\emph{The manuscript has undergone three rounds of revision and is now highly polished. The authors have addressed the reviewers' concerns regarding comparison with the latest state-of-the-art models, scalability for Big Data paradigms, experimental reproducibility, and terminology consistency. The paper provides a meaningful incremental advance in margin distribution learning. While the $O(n^3)$ complexity is a hurdle for ``Big Data,'' the provided theoretical paths toward distributed and approximated versions justify its publication.}
}


{\color{customgreen}
\subsubsection*{\underline{\textbf{Reviewer 7, Final Suggestion}}}
\emph{Before final publication, ensure that the ``Supplementary Material'' links are finalized and that all magenta/orange highlighting from the revision process is removed for the final production version.}
}


{\color{black}
\subsubsection*{\textbf{\underline{Author response}}}
}

{\color{black}
\subsubsection*{\textbf{\underline{Author action}}}
}

\newpage

% ---------------- Reviewer 8 ---------------------

{\color{black}
\section*{\textbf{REVIEWER 8}}\label{rev8}}

{\color{black}
\subsubsection*{\underline{General Comments}}
\emph{The authors have moved the manuscript forward in this round. The new Remark 2 clarifying ``quantile distance'' versus ``composite distance'' is useful, the MLP+Focal configuration is now properly specified, and the limitations discussion in the conclusion is a real improvement. That said, the new head-to-head against DeepSMOTE, LDAM, and Balanced Softmax introduces problems that need to be resolved, and several issues that should have been caught in earlier rounds are still present. I recommend major revision.}
}


{\color{customgreen}
\subsubsection*{\underline{\textbf{Reviewer 8, Strength 1}}}
\emph{The new Remark 2 clearly distinguishes the two distance concepts that were previously conflated. This was a real source of confusion in R2 and the fix is well done.}
}


{\color{customgreen}
\subsubsection*{\underline{\textbf{Reviewer 8, Strength 2}}}
\emph{The expanded MLP+Focal configuration in Section 5.3 (architecture, optimizer, $\beta$ values, batch size, LR schedule, grid search ranges) is detailed enough for reproducibility.}
}


{\color{customgreen}
\subsubsection*{\underline{\textbf{Reviewer 8, Strength 3}}}
\emph{The limitations paragraph in the conclusion (training complexity, extreme imbalance, concept drift, unstructured data) gives a more balanced positioning of the method than earlier versions.}
}


{\color{customgreen}
\subsubsection*{\underline{\textbf{Reviewer 8, Strength 4}}}
\emph{The scalability discussion (ADMM, mini-batch, Random Fourier Features) at least acknowledges the journal's scope, even if it is not yet demonstrated.}
}


{\color{customgreen}
\subsubsection*{\underline{\textbf{Reviewer 8, Major Concern 1}}}
\emph{The LDAM configuration as described is methodologically incoherent. The authors write that LDAM ``adapts ResNet-18 for tabular data.'' ResNet-18 is a 2D convolutional architecture designed for image inputs. There is no standard way to apply it to tabular features of dimension 3 to 61, which is the range in Table 1. The authors need to explain exactly what architecture was used, replace it with a sensible tabular backbone (MLP with LDAM loss would be the natural choice), or drop the claim. As it stands, readers cannot reproduce or evaluate this baseline.}
}


{\color{black}
\subsubsection*{\textbf{\underline{Author response}}}
}


{\color{black}
\subsubsection*{\textbf{\underline{Author action}}}
}


{\color{customgreen}
\subsubsection*{\underline{\textbf{Reviewer 8, Major Concern 2}}}
\emph{Please release the code, seeds, and fold-level outputs for Table 7. The ordering LDAM over Balanced Softmax over DeepSMOTE is very consistent across the 17 datasets and across Accuracy, AUC, and F1. This ordering is plausible on its own, but the consistency is unusual enough that the paper should provide a reproducibility package (code, random seeds, per-fold numbers). This would also address Reviewer 5's original reproducibility request in a more complete way.}
}


{\color{black}
\subsubsection*{\textbf{\underline{Author response}}}
}


{\color{black}
\subsubsection*{\textbf{\underline{Author action}}}
}


{\color{customgreen}
\subsubsection*{\underline{\textbf{Reviewer 8, Major Concern 3}}}
\emph{The Wilcoxon p-value reporting needs more detail. Reporting p = 0.03125 for all four comparisons (DBUPLDM vs SVM, UPSVM, PinSVM, LDM) is consistent with a test over 5 or 6 paired observations where DBUPLDM wins all of them, but the paper does not say what the paired units are (datasets? $\eta$ values? $\eta \times$ datasets?) or what n is. Please state the test setup and report effect sizes alongside the p-values. A one-line description would clear this up.}
}


{\color{black}
\subsubsection*{\textbf{\underline{Author response}}}
}


{\color{black}
\subsubsection*{\textbf{\underline{Author action}}}
}



{\color{customgreen}
\subsubsection*{\underline{\textbf{Reviewer 8, Major Concern 4}}}
\emph{Text-table inconsistency. Section 5.3 says ``deep learning methods show strong performance on certain datasets (e.g., LDAM on Ionosphere and Australia).'' In Table 7, DBUPLDM beats LDAM on Australian (87.24 vs 84.20). Pick examples that actually match the table, or remove the specific claim.}
}


{\color{black}
\subsubsection*{\textbf{\underline{Author response}}}
}


{\color{black}
\subsubsection*{\textbf{\underline{Author action}}}
}


{\color{customgreen}
\subsubsection*{\underline{\textbf{Reviewer 8, Major Concern 5}}}
\emph{Big Data Research scope is still only rhetorical. The new scalability paragraph in the conclusion discusses ADMM, mini-batch processing, and Random Fourier Features, but none of this is implemented or evaluated. The experiments are entirely on small UCI datasets (hundreds to about a thousand samples). At minimum, I would ask for one empirical demonstration on a larger-scale imbalanced dataset (for example, CreditCard Fraud with about 285k samples, KDD99, or a similar public benchmark) to support the scalability claims. Without this, the journal fit remains weak.}
}


{\color{black}
\subsubsection*{\textbf{\underline{Author response}}}
}


{\color{black}
\subsubsection*{\textbf{\underline{Author action}}}
}


{\color{customgreen}
\subsubsection*{\underline{\textbf{Reviewer 8, Major Concern 6}}}
\emph{Novelty over the authors' own prior work should be spelled out. The UP loss applied to a large-margin distribution classifier already appears in Gu et al. 2024 (cited in the paper). The dual balance factor is close in spirit to Cheng et al.'s CS-LDM (2016). The paper's contribution is essentially the combination of these two ingredients. A short paragraph in the introduction or Section 3 stating what is new here relative to Gu et al. 2024 would strengthen the case.}
}


{\color{black}
\subsubsection*{\textbf{\underline{Author response}}}
}


{\color{black}
\subsubsection*{\textbf{\underline{Author action}}}
}


{\color{customgreen}
\subsubsection*{\underline{\textbf{Reviewer 8, Minor Issue 1}}}
\emph{Title: ``Dual Balanced Large Marginal Machine'' should be ``Large Margin Machine'' (or ``Large Margin Distribution Machine'' to match the abstract). This has remained wrong across rounds and needs to be fixed.}
}


{\color{black}
\subsubsection*{\textbf{\underline{Author response}}}
}


{\color{black}
\subsubsection*{\textbf{\underline{Author action}}}
}


{\color{customgreen}
\subsubsection*{\underline{\textbf{Reviewer 8, Minor Issue 2}}}
\emph{Dataset naming is inconsistent. Table 1 lists ``Herberman,'' the results tables use ``Heberman,'' and Table 7 uses ``Habe.'' The dataset is Haberman's Survival. Pick one spelling and use it throughout.}
}


{\color{black}
\subsubsection*{\textbf{\underline{Author response}}}
}


{\color{black}
\subsubsection*{\textbf{\underline{Author action}}}
}


{\color{customgreen}
\subsubsection*{\underline{\textbf{Reviewer 8, Minor Issue 3}}}
\emph{``Ecoil'' should be ``Ecoli'' (E. coli dataset).}
}


{\color{black}
\subsubsection*{\textbf{\underline{Author response}}}
}


{\color{black}
\subsubsection*{\textbf{\underline{Author action}}}
}


{\color{customgreen}
\subsubsection*{\underline{\textbf{Reviewer 8, Minor Issue 4}}}
\emph{``Germans'' should likely be ``German'' (German Credit dataset).}
}


{\color{black}
\subsubsection*{\textbf{\underline{Author response}}}
}


{\color{black}
\subsubsection*{\textbf{\underline{Author action}}}
}


{\color{customgreen}
\subsubsection*{\underline{\textbf{Reviewer 8, Minor Issue 5}}}
\emph{``Phonemes'' vs ``Phoneme'' is inconsistent between Table 1 and Section 5.5.}
}


{\color{black}
\subsubsection*{\textbf{\underline{Author response}}}
}


{\color{black}
\subsubsection*{\textbf{\underline{Author action}}}
}


{\color{customgreen}
\subsubsection*{\underline{\textbf{Reviewer 8, Minor Issue 6}}}
\emph{Table 7, WDBC row: ``87.50/91.43/0/89'' has a slash where a decimal should be. Should read ``0.89.''}
}


{\color{black}
\subsubsection*{\textbf{\underline{Author response}}}
}


{\color{black}
\subsubsection*{\textbf{\underline{Author action}}}
}


{\color{customgreen}
\subsubsection*{\underline{\textbf{Reviewer 8, Minor Issue 7}}}
\emph{Figure 6, subfigures (i), (j), (k), and (l) are all captioned ``Prlx dataset.'' These are clearly different datasets. Fix the captions.}
}


{\color{black}
\subsubsection*{\textbf{\underline{Author response}}}
}


{\color{black}
\subsubsection*{\textbf{\underline{Author action}}}
}


{\color{customgreen}
\subsubsection*{\underline{\textbf{Reviewer 8, Minor Issue 8}}}
\emph{Figure 6(g) caption reads ``Daibetes dataset.'' Should be ``Diabetes.''}
}


{\color{black}
\subsubsection*{\textbf{\underline{Author response}}}
}


{\color{black}
\subsubsection*{\textbf{\underline{Author action}}}
}


{\color{customgreen}
\subsubsection*{\underline{\textbf{Reviewer 8, Minor Issue 9}}}
\emph{Section 5.5: ``PinSVM, UPSVM, and DBUPLDM) performe well'' should be ``perform.''}
}


{\color{black}
\subsubsection*{\textbf{\underline{Author response}}}
}


{\color{black}
\subsubsection*{\textbf{\underline{Author action}}}
}


{\color{customgreen}
\subsubsection*{\underline{\textbf{Reviewer 8, Minor Issue 10}}}
\emph{The Section 1 roadmap reads ``Section 3 proposes the Dual Balanced method and DBUPLDM model. Section 5 discusses...'' Section 4 (Properties of DBUPLDM) is skipped. Add it.}
}


{\color{black}
\subsubsection*{\textbf{\underline{Author response}}}
}


{\color{black}
\subsubsection*{\textbf{\underline{Author action}}}
}


{\color{customgreen}
\subsubsection*{\underline{\textbf{Reviewer 8, Minor Issue 11}}}
\emph{Table 5 caption ``Minimum distance with second order statistics'' does not describe the content (ablation over balanced factors and hinge vs UP loss). Rewrite.}
}


{\color{black}
\subsubsection*{\textbf{\underline{Author response}}}
}


{\color{black}
\subsubsection*{\textbf{\underline{Author action}}}
}


{\color{customgreen}
\subsubsection*{\underline{\textbf{Reviewer 8, Minor Issue 12}}}
\emph{Remark renumbering: the new Remark 2 (Terminology Clarification) shifts former Remark 2 to Remark 3 and former Remark 3 to Remark 4. Please verify no downstream cross-references still point to the old numbers.}
}


{\color{black}
\subsubsection*{\textbf{\underline{Author response}}}
}


{\color{black}
\subsubsection*{\textbf{\underline{Author action}}}
}


{\color{customgreen}
\subsubsection*{\underline{\textbf{Reviewer 8, Minor Issue 13}}}
\emph{The third contribution bullet in the introduction says ``the quantile distance is introduced to replace the shortest distance between sample sets in traditional LDM.'' Given the new Remark 2, this is imprecise (the replacement involves both the UP loss and the composite distance). Reword for consistency.}
}


{\color{black}
\subsubsection*{\textbf{\underline{Author response}}}
}


{\color{black}
\subsubsection*{\textbf{\underline{Author action}}}
}


{\color{customgreen}
\subsubsection*{\underline{\textbf{Reviewer 8, Minor Issue 14}}}
\emph{The response letter header says ``BDR-D-25-00303R3'' but the opening paragraph says ``R1,'' and later describes this as the third round. Fix this.}
}


{\color{black}
\subsubsection*{\textbf{\underline{Author response}}}
}


{\color{black}
\subsubsection*{\textbf{\underline{Author action}}}
}


{\color{customgreen}
\subsubsection*{\underline{\textbf{Reviewer 8, Minor Issue 15}}}
\emph{Reviewer 6's comment about ``deeper theoretical insight'' is acknowledged in the response but not concretely addressed. A brief sentence on what was done, or why no change was needed, would help.}
}


{\color{black}
\subsubsection*{\textbf{\underline{Author response}}}
}


{\color{black}
\subsubsection*{\textbf{\underline{Author action}}}
}


\newpage

\newpage
\section*{\textbf{ACKNOWLEDGEMENT}}
\noindent\rule[0.25\baselineskip]{\textwidth}{1pt}
% {\color{black}
% \section*{\textbf{ACKNOWLEDGEMENT}}}
% 特别感谢编辑和匿名评论者的宝贵和有益的意见和建议。

Special thanks to you for the valuable and helpful comments and suggestions from editors and anonymous reviewers. 
% 这些评论有助于改进我们的论文，也对我们的研究具有重要的指导意义。
The comments are helpful in improving our paper, as well as the important guiding significance to our research. 
% 我们希望编辑和审稿人的关切在修订版中得到充分解决。
We hope that the editors' and reviewers' concerns have been adequately addressed in the revised version.

Yours sincerely,

\rightline{Pingping Gu, Denghao Dong, Quan Shi, Fudi Wang and Guojun Huang*}



%\bibliography{IEEEabrv,mybib}
%\bibliographystyle{IEEEtran}
%\bibliographystyle{ieeetr}

\end{document}
